\chapter*{Extended Abstract}
\addcontentsline{toc}{chapter}{Extended Abstract}
\markboth{Extended Abstract}{Extended Abstract}

\section*{Context and objective}

Warehouses serving e-commerce demand face a monthly planning decision with two components that
are usually treated separately: how much labour capacity to assign to each stage of the
fulfilment process, and in what order to release the work that queues at those stages. Capacity
is typically sized from aggregate workload ratios, while the sequencing rule is treated as a
fixed operating condition. Two features of modern fulfilment make that separation
unsatisfactory. Demand is bursty and seasonal, so queueing effects that aggregate ratios cannot
represent determine whether a plan actually holds. And service commitments are differentiated:
urgent and normal orders carry very different deadlines, so a single aggregate service figure
can conceal complete failure of the more demanding class.

The objective of this thesis is to design, implement and evaluate a simulation-based
decision-support framework that recommends a monthly workforce configuration and an order
sequencing policy jointly, under class-specific service constraints and an explicit economic
objective, and that explains its recommendations well enough to be acted on.

\section*{Theoretical background}

The work draws on four bodies of literature. Warehouse operations research establishes order
picking as the dominant component of warehouse operating cost. Queueing theory explains why
capacity sized to average requirements does not deliver average performance, since waiting grows
non-linearly with utilisation. Scheduling research establishes that rules using due-date
information outperform first-in-first-out under congestion, while priority queueing establishes
that a non-pre-emptive priority discipline redistributes waiting between classes rather than
removing it. Reinforcement learning provides a means of deriving a sequencing policy from
experience rather than specifying one in advance.

The position this thesis occupies is one of integration rather than novelty in any single
component. Existing work generally treats capacity and sequencing as separate problems and
evaluates sequencing rules against aggregate objectives; this study treats them as one joint
decision evaluated against a conjunction of class-specific constraints.

\section*{Methodology}

A discrete-event simulation model represents a three-stage fulfilment process --- picking,
packing and dispatch --- each with its own queue and dedicated worker pool. Orders are
heterogeneous, and critically their workload is stage-differentiated: product family,
complexity and urgency multiply an order's workload differently at each stage, so a fragile
order is slightly harder to pick but substantially harder to pack.

Three design decisions distinguish the model. First, the simulated physical capacity and the
economically paid capacity are the same quantity by construction: the monthly horizon is
defined as the paid hours of one full-time equivalent, so a worker can never process work it is
not paid for. Second, the horizon is finite and work left unfinished at its end is retained as
explicit backlog counting as service failure, which prevents insufficient capacity from being
concealed by unlimited post-period processing. Third, feasibility is conjunctive: a
configuration is acceptable only if it meets both the urgent and the normal service floors
separately.

Three sequencing policies are compared: first-in-first-out; a strict two-class priority rule;
and a deep Q-network agent, shared across all three stages, that chooses between an urgent and a
normal queue using a sixteen-feature state and receives a reward deferred until dispatch. The
agent observes only present state and has no information about future arrivals. All three are
compared under common random numbers --- service times are pre-computed per order and stage
before any policy runs --- so differences reflect only the sequencing decision.

An analytical capacity estimate anchors a bounded search over sixteen structured workforce
candidates per month, each simulated under each policy. A transparent weighted diagnostic ranks
the stages by operational pressure, and an adaptive step tests whether adding one worker at the
identified bottleneck is economically justified. The framework was exercised retrospectively
over twelve months of order-level data and prospectively over replicated scenarios for a peak
month. All data are synthetic. The retrospective experiments use a synthetic annual dataset of
\num{240000} order-level records; the forecast experiment does not replay that dataset but
generates fresh stochastic order streams from an aggregate December volume assumption and the
configured demand profile. No real, proprietary or client data were used.

\section*{Main findings}

The analytical estimate exceeded the simulated recommendation in \emph{all twelve} months, by up
to nine full-time equivalents in December, where the simulated recommendation of 38 FTE lay
approximately nineteen per cent below the analytical estimate of 47. Two separate mechanisms
produce this: independent per-stage rounding, which persistently over-provides the smaller
stages, and a substitution between capacity and operating discipline that no aggregate formula
can represent.

First-in-first-out satisfied both service floors in only 84 of 192 configurations, against 166
and 167 for the two urgency-aware policies. More decisively, its \emph{complete} monthly
infeasibility is confined to the four campaign months: in January, February, November and
December it met both floors at none of the sixteen candidate workforces. Within the workforce
range evaluated, therefore, additional capacity did not make first-in-first-out feasible in
those months; a sufficiently large workforce would eventually do so, but no candidate examined
reached it.

The value of sequencing proved strongly conditional on the capacity regime. Where capacity was
ample the three policies were practically indistinguishable, differing in total cost by thirty
euros at one October configuration. Where capacity was tight they diverged sharply. At
December's recommended configuration the strict priority rule protected the urgent class but
drove the normal class below its floor and was infeasible, while the learned policy held both
classes above their floors at thirteen per cent lower cost. Expressed in capacity terms, the
strict priority rule required seven more full-time equivalents than the learned policy to reach
feasibility at all in that month. These are differences between the policies as implemented: the
two engines differ both in how the serving class is chosen and in how orders are ordered within
a class, so the gap cannot be attributed to the learned decision rule alone.

Because policies were compared at identical workforces, labour cost was constant and the entire
economic difference lay in late-order penalties --- negligible for most of the year and dominant
in the peak. Backlog and utilisation were nearly invariant across policies, confirming that
sequencing reallocates lateness rather than creating capacity.

Picking was identified as the primary pressure location in all twelve months, accounting for
between eighty-four and one hundred per cent of the waiting accumulated by late orders. Yet in
nine months the framework tested adding one worker there and rejected it every time: a worker
costing \num{2400} euros per month avoided under \num{130} euros of penalties in six of them.

\section*{Discussion, contributions and conclusion}

The results support a single organising conclusion: capacity and sequencing are not separable
decisions, and treating them separately is expensive in a measurable way. They also yield two
findings that run against intuition --- that a correctly identified bottleneck is frequently the
wrong place to add capacity, and that a learned policy can be genuinely valuable in a narrow
regime while being indistinguishable from a one-line heuristic for most of the year.

The contributions are methodological (an operating-time design equating physical and paid
capacity, a finite-horizon formulation with explicit backlog, conjunctive class-specific
feasibility, and a policy comparison under common random numbers), practical (a workflow
producing an interrogable recommendation with a named bottleneck and an economic test of
marginal capacity), and empirical (evidence that sequencing value is capacity-regime dependent).

The principal limitations are that all evidence is synthetic and uncalibrated, that the model
omits warehouse physics such as travel and layout, that the retrospective results rest on one
demand realisation per month, that the two compared engines differ in their within-class queue
mechanics, and that the trained policy's weights were not retained and so its results could not
be regenerated. The framework is therefore a coherent instrument for comparing plans under
stated assumptions, and is not established as predictive of any particular operation.
